\documentclass[11pt]{article}

\usepackage[utf8]{inputenc}
\usepackage[T1]{fontenc}
\usepackage{lmodern}
\usepackage{geometry}
\usepackage{amsmath,amssymb,amsfonts,amsthm,mathtools}
\usepackage{enumitem}
\usepackage{microtype}
\usepackage{hyperref}

\geometry{margin=1in}
\hypersetup{colorlinks=true, linkcolor=black, urlcolor=blue, citecolor=black}
\setlist[enumerate]{topsep=0.35em,itemsep=0.3em}
\setlist[itemize]{topsep=0.3em,itemsep=0.25em}

\newcommand{\Aether}{\AE{}ther}
\newcommand{\Aflow}{\AE{}ther-flow}
\newcommand{\TheoryProgram}{\Aether{} / \Aflow{} framework}
\newcommand{\Sub}{\mathcal{A}}
\newcommand{\Step}{\mathsf{Step}}
\newcommand{\RespAlg}{\mathcal{R}_{\Gamma}}
\newcommand{\Rank}{\mathsf{rank}_{\mathcal{S}}}
\newcommand{\Cycle}{\mathcal{C}_{\Gamma}}
\newcommand{\Lang}{\mathcal{L}_{0}}
\newcommand{\Sel}{\mathfrak{S}}

\newtheorem{auditfinding}{Audit finding}
\newtheorem{smugglingrisk}{Smuggling risk}
\newtheorem{promotionblocker}{Promotion blocker}
\newtheorem{auditverdict}{Audit verdict}
\newtheorem{nextburden}{Next burden}

\title{\textbf{Localization Source-Basis Axiom Selection-Law Smuggling Audit}\\[0.35em]
\large Target-Import Review of Source-Internal Selector Rules}
\author{Smuggling Auditor AgentJob AJ-RT-20260608-014-001}
\date{June 9, 2026}

\begin{document}

\maketitle

\begin{abstract}
This draft audits the source-internal selection-law packet
\(\Sel=(S_0,S_1,S_2,S_3,S_4,S_F)\) for hidden target imports. The packet is a
useful control improvement: it makes selection rules explicit, requires
benchmark blindness, and preserves underselection as a blocker. The audit
nevertheless finds that the selectors are not promotion-ready. Presentation
seeds, refinement policy, transition minimality, discriminator-family
closure, quotient admissibility, shared-support transport, and
failure-status semantics remain possible target-import channels. Candidate
reconstruction and Gate Chair review remain blocked.
\end{abstract}

\section{Scope and Audit Standard}

This artifact is a draft/control output for task \texttt{RT-20260608-014}. It
does not edit canonical ontology TeX, adopt the selection laws as canonical
ontology, authorize candidate reconstruction, issue a Gate Chair verdict, or
reject the broader \TheoryProgram{}.

The audit asks two questions:
\begin{enumerate}
    \item Does the packet directly introduce target metric, target null-cone,
    target coordinate, target observer protocol, target rank, or target
    proper-time data as source input?
    \item Even if it avoids direct import, can its source-named selection
    rules be chosen or interpreted so that target structures enter before
    benchmark comparison?
\end{enumerate}

The first question is a declaration screen. The second question is the
promotion-readiness screen.

\section{Positive Transparency Screen}

\begin{auditfinding}[Explicit selector surface]
The packet correctly exposes the selection surface. It names \(S_0\)--\(S_4\)
and \(S_F\), requires source-side definition before benchmark comparison, and
requires non-invariant outputs to be reported instead of silently repaired.
\end{auditfinding}

\begin{auditfinding}[No direct target primitive]
The packet does not explicitly install a Lorentzian metric, continuum
manifold topology, target null cone, radar-coordinate protocol, target
proper-time normalization, or fixed rank four as a source primitive.
\end{auditfinding}

These are control successes. They do not prove that the selectors are
source-internally forced, and they do not authorize candidate reconstruction.

\section{S0: Presentation Seed and Refinement Risks}

\begin{smugglingrisk}[Source-accessible seed can select target-like locality]
\(S_0\) uses a bounded source-accessibility seed and a refinement policy
\(\rho_U\). The packet forbids choosing these because they resemble target
neighborhoods, but it does not yet supply an independent source law that
fixes the seed, closure scale, or refinement policy. A benchmark-compatible
seed can still select a target-like localization arena before the audit can
observe the target import.
\end{smugglingrisk}

\begin{promotionblocker}[Presentation-source blocker]
\(S_0\) is not promotion-ready until the seed and refinement policy are fixed
by source dynamics or until a Refuter shows that all admissible seeds and
refinements preserve the same localization-relevant outputs.
\end{promotionblocker}

\section{S1: Transition Minimality Risks}

\begin{smugglingrisk}[Minimal support can hide topology]
\(S_1\) selects \(\Step\) by order comparability, shared support, and absence
of an intermediate record with the same support class. The terms
``shared support'' and ``same support class'' remain doing substantive work.
If those classes are selected to recover target-like fronts or adjacency
balls, target topology has entered through the support equivalence rather than
through an explicit metric.
\end{smugglingrisk}

\begin{promotionblocker}[Transition-family blocker]
The transition selector can support later construction only after the
admissible support equivalence is source-fixed, or after all admissible
transition families are surveyed and non-invariant fronts are preserved as
negative results.
\end{promotionblocker}

\section{S2: Discriminator Closure Risks}

\begin{smugglingrisk}[All admissible families can still be pre-shaped]
\(S_2\) improves the prior state by carrying all source-defined perturbation
and discriminator families rather than a hand-selected member. The remaining
risk is the definition of the family itself. Recurrence words, transition
windows, and response distinctions can be made broad or narrow enough to
recover benchmark-compatible first-response ranks.
\end{smugglingrisk}

\begin{promotionblocker}[Family-definition blocker]
\(\mathcal{P}\) and \(\mathcal{D}\) remain blocked until their closure rules
are fixed before benchmark comparison and shown not to vary with benchmark
success. Disagreement among allowed discriminators must remain a blocker.
\end{promotionblocker}

\section{S3: Quotient and Rank Risks}

\begin{smugglingrisk}[Rank can enter through quotient admissibility]
\(S_3\) states that rank four is not imposed. That is necessary. It is not
sufficient. The admissible response-word ground set, recurrence equivalence,
and recombination rules can still be chosen so that the quotient produces the
desired response rank.
\end{smugglingrisk}

\begin{promotionblocker}[Rank-spectrum blocker]
A later candidate may not use a favorable rank unless the rank is invariant
across source-admissible quotient choices. Variable rank is a result, not a
defect to be tuned away.
\end{promotionblocker}

\section{S4: Clock-Transport Risks}

\begin{smugglingrisk}[Shared-support transport can encode synchronization]
\(S_4\) requires transport from shared source transition support and requires
holonomy reporting. The remaining risk is that shared support can encode an
observer-synchronization convention, connection-like rule, or proper-time
scale while remaining source-named. Reporting holonomy is necessary, but it
does not by itself prove that the transport law is source-internal.
\end{smugglingrisk}

\begin{promotionblocker}[Transport-source blocker]
Clock comparison remains blocked until the shared-support relation and
transport rule are fixed by source dynamics, or until a failure-tolerant
survey shows that transported cycle ratios are invariant across admissible
transport choices.
\end{promotionblocker}

\section{S_F: Failure-Status Risks}

\begin{smugglingrisk}[Status semantics can defer the burden]
\(S_F\) is a useful safeguard because only \(\mathsf{invariant}\) status may
support later candidate objects. The risk is semantic drift. A
\(\mathsf{not\_computed}\) or \(\mathsf{family\_valued}\) status must not be
treated as partial success, and \(\mathsf{invariant}\) must not be asserted
without an explicit survey domain.
\end{smugglingrisk}

\begin{promotionblocker}[Failure-ledger blocker]
The failure ledger must name the admissible class surveyed, the source inputs
used, and the benchmark inputs excluded. Without that metadata, the status
labels are bookkeeping rather than evidence.
\end{promotionblocker}

\section{Audit Verdict}

\begin{auditverdict}[Transparent selectors, unresolved target-import channels]
The selection-law packet passes a narrow declaration audit: it makes selector
rules explicit and does not directly install target metric, target null-cone,
target coordinates, target rank, or target proper-time normalization. It
fails promotion-readiness. The unresolved target-import channels are:
\begin{enumerate}
    \item source-accessible seeds and refinement policy can select target-like
    local arenas;
    \item support equivalence in \(S_1\) can encode target topology through
    transition incidence;
    \item discriminator-family closure in \(S_2\) can tune first-response
    ranks;
    \item quotient admissibility in \(S_3\) can manufacture favorable response
    rank; and
    \item shared-support transport in \(S_4\) can encode observer
    synchronization or clock normalization.
\end{enumerate}
\end{auditverdict}

This is a local control result. It does not reject the \TheoryProgram{} and
does not reject the use of explicit selection laws as research objects. It
does block candidate reconstruction from the current packet.

\section{Next Burden}

\begin{nextburden}[Refuter benchmark-independence test]
The logical next role is a Refuter AgentJob that tests whether the selection
laws are benchmark-independent over their admissible source-defined classes.
The test should require explicit survey domains for seeds, refinements,
support equivalences, discriminator families, response quotients, and
transport rules. If any localization-relevant output is non-invariant, the
result should remain a negative control result or route to repair, not to
candidate reconstruction.
\end{nextburden}

Candidate Constructor, canonical ontology edit, and Gate Chair review remain
premature.

\section*{References}

Ricciardi, A. S. (2026, June 9). \emph{Localization source-basis axiom
selection-law packet: Source-internal selector rules, invariance obligations,
and failure reporting for A0--A4} [Unpublished manuscript]. The Aether-Flow
Interpretation of Relativity Research Project,
\texttt{research\_control/tasks/RT-20260608-013/artifacts/12\_LOCALIZATION\_SOURCE\_BASIS\_AXIOM\_SELECTION\_LAW\_PACKET.tex}.

Ricciardi, A. S. (2026, June 9). \emph{Localization source-basis axiom
benchmark-independence refutation: Source-internal necessity test for A1--A4}
[Unpublished manuscript]. The Aether-Flow Interpretation of Relativity
Research Project,
\texttt{research\_control/tasks/RT-20260608-012/artifacts/11\_LOCALIZATION\_SOURCE\_BASIS\_AXIOM\_BENCHMARK\_INDEPENDENCE\_REFUTATION.tex}.

Ricciardi, A. S. (2026, June 9). \emph{Localization source-basis axiom
smuggling audit: Target-import review of A1--A4 extension hypotheses}
[Unpublished manuscript]. The Aether-Flow Interpretation of Relativity
Research Project,
\texttt{research\_control/tasks/RT-20260608-011/artifacts/10\_LOCALIZATION\_SOURCE\_BASIS\_AXIOM\_SMUGGLING\_AUDIT.tex}.

\end{document}
